\section{Présentation générale}

\subsection{Objectif}

Cette application web permet l'apprentissage du vocabulaire ukrainien à travers une interface interactive. Elle présente des mots classés par catégorie grammaticale (noms, verbes, adjectifs, etc.) avec leurs déclinaisons ou conjugaisons complètes, des accents toniques visuels, et des phrases d'exemple.

L'application est construite avec \textbf{SvelteKit} (framework web basé sur Svelte~5) et déployée en site statique via \texttt{adapter-static}.

\subsection{Fonctionnalités principales}

\begin{enumerate}
  \item \textbf{Navigation par catégorie} --- La sidebar gauche liste tous les mots, regroupés par catégorie grammaticale (\ukr{nom}, \ukr{adj.}, \ukr{verbe}, etc.).

  \item \textbf{Fiches détaillées} --- Un clic sur un mot affiche ses formes fléchies dans le panneau central : tableau de déclinaison (noms, adjectifs) ou de conjugaison (verbes).

  \item \textbf{Accents toniques} --- Une case à cocher (en bas à droite) active ou désactive l'affichage des accents. Par exemple, \ukr{оди́н} avec accent vs \ukr{один} sans accent.

  \item \textbf{Info-bulle au survol} --- Le survol d'un mot ukrainien (span \texttt{.ukr}) affiche une bulle avec le lemme accentué et les métadonnées grammaticales (catégorie, cas, nombre, genre).

  \item \textbf{Tableau de grammaire épinglable} --- Un clic sur un mot ukrainien «~épingle~» la sidebar droite, qui affiche le tableau complet de déclinaison/conjugaison. La forme survolée est mise en gras.

  \item \textbf{Couleur par cas grammatical} --- Au survol, la couleur du mot reflète son cas : nominatif (hérité), génitif (vert), accusatif (bleu), locatif (brun), instrumental/datif (rose), vocatif (rouge foncé).

  \item \textbf{Phrases d'exemple} --- Chaque mot peut avoir des phrases associées, affichées sous la fiche détaillée avec traduction française.

  \item \textbf{Recherche de phrases} --- Une page dédiée (\texttt{/phrases}) permet de filtrer les phrases par métadonnées (auteur, cours, source).

  \item \textbf{Génération de conjugaisons} --- Dans les phrases d'exemple, une case à cocher déploie toutes les formes conjuguées d'un verbe pour un temps donné.

  \item \textbf{Responsive} --- L'interface s'adapte aux écrans étroits (empilement des colonnes sous 768px).
\end{enumerate}

\subsection{Captures d'écran}

% ============ PLACEHOLDER ============
% Capture : page principale complète avec la sidebar gauche (liste des mots)
% et le panneau central (détails d'un nom, par exemple "будинок")
% Taille recommandée : pleine largeur du navigateur, ~1200x800px
\begin{figure}[ht]
  \centering
  \fbox{\parbox{0.8\textwidth}{\centering\vspace{3cm}\textit{[screenshot-main.png]\\Page principale : sidebar + détails}\vspace{3cm}}}
  % \includegraphics[width=0.9\textwidth]{screenshot-main.png}
  \caption{Page principale --- liste des mots (gauche) et fiche détaillée (centre).}
  \label{fig:main}
\end{figure}

% ============ PLACEHOLDER ============
% Capture : page /phrases avec la barre de recherche en haut à droite
% et quelques phrases affichées avec leur traduction
\begin{figure}[ht]
  \centering
  \fbox{\parbox{0.8\textwidth}{\centering\vspace{3cm}\textit{[screenshot-phrases.png]\\Page de recherche de phrases}\vspace{3cm}}}
  % \includegraphics[width=0.9\textwidth]{screenshot-phrases.png}
  \caption{Page de recherche de phrases avec filtre par métadonnées.}
  \label{fig:phrases}
\end{figure}

\subsection{Installation et lancement}

\begin{lstlisting}[language=bash, title=Commandes de base]
npm install          # Installer les dependances
npm run dev          # Serveur de developpement (hot reload)
npm run build        # Generer le site statique dans build/
npm run preview      # Previsualiser le site genere
npm run test         # Lancer les tests unitaires (Vitest)
\end{lstlisting}
